\documentclass[12pt,a4paper]{article}
\usepackage[utf8]{inputenc}
\usepackage[T1]{fontenc}
\usepackage[polish]{babel}
\usepackage{geometry}
\usepackage{enumitem}
\usepackage{titlesec}
\usepackage{parskip}
\usepackage{fancyhdr}
\usepackage{xcolor}
\usepackage{mdframed}
\usepackage{microtype}
\usepackage{booktabs}
\usepackage{hyperref}
\usepackage{multirow}
\usepackage{array}
\usepackage{longtable}
\usepackage{tabularx}

\geometry{margin=2.5cm}
\hypersetup{colorlinks=true, linkcolor=black, urlcolor=blue}

\pagestyle{fancy}
\fancyhf{}
\rhead{SDLC-LLM Benchmark}
\lhead{Zadanie R2 -- Refinement (Hybrid)}
\rfoot{\thepage}

\titleformat{\section}{\large\bfseries}{}{0em}{}[\titlerule]
\titleformat{\subsection}{\normalsize\bfseries}{}{0em}{}
\titleformat{\subsubsection}{\normalsize\itshape}{}{0em}{}

\definecolor{metabg}{gray}{0.93}
\definecolor{promptbg}{RGB}{235,245,255}
\definecolor{claudecolor}{RGB}{210,105,30}
\definecolor{score3}{RGB}{220,240,220}
\definecolor{score2}{RGB}{255,243,205}
\definecolor{score1}{RGB}{255,220,200}
\definecolor{score0}{RGB}{255,200,200}
\definecolor{tpcolor}{RGB}{220,240,220}
\definecolor{fpcolor}{RGB}{255,220,200}
\definecolor{fncolor}{RGB}{255,243,205}

\newcommand{\scorebox}[1]{%
  \ifnum#1=3\colorbox{score3}{\textbf{3}}%
  \else\ifnum#1=2\colorbox{score2}{\textbf{2}}%
  \else\ifnum#1=1\colorbox{score1}{\textbf{1}}%
  \else\colorbox{score0}{\textbf{0}}\fi\fi\fi
}
\newcommand{\tp}[1]{\colorbox{tpcolor}{\small #1}}
\newcommand{\fp}[1]{\colorbox{fpcolor}{\small #1}}
\newcommand{\fn}[1]{\colorbox{fncolor}{\small #1}}

\begin{document}

% ─── STRONA TYTUŁOWA ────────────────────────────────────────────────
\begin{center}
  {\LARGE\bfseries Requirements Artifact}\\[1.5em]
  {\large \textbf{Case study:} Appointment Booking System}\\[0.3em]
  {\large \textbf{Model:} \texttt{claude-sonnet-4-6}}\\[0.3em]
  {\large \textbf{Tryb:} \textbf{Hybrid} (LLM + Human-in-the-loop)}\\[0.3em]
  {\large \textbf{Wejście:} Gold ver2 (challenging -- celowe niejednoznaczności)}
\end{center}

\vspace{1em}\hrule\vspace{1.5em}

% ─── WPROWADZONE ZMIANY (HUMAN CHANGES) ─────────────────────────────
\section{Interwencja analityka (Human Changes)}

\begin{mdframed}[backgroundcolor=promptbg, linewidth=0pt, innerleftmargin=10pt, innerrightmargin=10pt, innertopmargin=8pt, innerbottommargin=8pt]
\small
\textbf{Changelog z przeglądu manualnego:}\\
FR4: usunięto ogólne sformułowanie 'w dozwolonym oknie czasowym', dodano sztywny limit 'do 24 godzin przed terminem'. 
FR5: usunięto możliwość modyfikacji terminów dla już aktywnych rezerwacji. 
FR6: zastąpiono stary punkt wymaganiem o historii i dodano dokładną listę statusów (AVAILABLE, BOOKED, CANCELLED, COMPLETED, BLOCKED). 
FR7: dodano wymóg ograniczający liczbę rezerwacji do 3 aktywnych na użytkownika. 
FR9: dodano wymóg obsługi równoczesnych prób rezerwacji (współbieżność). 
FR ogólne: usunięto stare FR6, FR7, FR8, FR12 (dotyczące powiadomień SMS/e-mail, zarządzania interfejsem ról oraz raportów administracyjnych), przenumerowano pozostałe. 
NFR4: usunięto wymóg uwierzytelniania MFA, dodano wymóg spójności stanu końcowego przy równoczesnych operacjach. 
NFR ogólne: usunięto NFR7, NFR9, NFR10, NFR11, NFR12 (dotyczące WCAG 2.1, SLA dla powiadomień, backupów, testów 80\% i OWASP), przenumerowano NFR8 na NFR7. 
US ogólne: usunięto US5, US6, US10, US13, US14, US15 (raporty i uprawnienia), skrócono listę z 15 do 10. 
AC ogólne: zachowano podstawowe przepływy (AC1-AC6), usunięto stare AC7-AC10, AC12-AC16, AC19-AC23 (kryteria dla powiadomień, weryfikacji historii, uprawnień i raportowania), redukując całkowitą liczbę AC z 23 do 9.
\end{mdframed}

% ─── CEL SYSTEMU ────────────────────────────────────────────────────
\section{Cel systemu}

System umożliwia rezerwację wizyt u specjalistów, zapewniając użytkownikom możliwość przeglądania dostępnych terminów, dokonywania i anulowania rezerwacji. Specjaliści zarządzają swoim grafikiem, a system zapobiega konfliktom rezerwacji przy zachowaniu skalowalności i obsłudze wielu ról użytkowników.

% ─── ROLE ───────────────────────────────────────────────────────────
\section{Role}

\begin{itemize}[leftmargin=1.5em]
  \item \textbf{Użytkownik} -- przeglądanie dostępnych terminów, dokonywanie rezerwacji wizyty, anulowanie własnej rezerwacji w dozwolonym oknie czasowym, przeglądanie historii własnych rezerwacji.
  \item \textbf{Specjalista} -- zarządzanie własnym grafikiem dostępności, modyfikowanie terminów istniejących rezerwacji, przeglądanie listy własnych rezerwacji, blokowanie wybranych terminów, anulowanie rezerwacji pacjenta z podaniem przyczyny.
  \item \textbf{Administrator} -- zarządzanie kontami użytkowników i specjalistów, przeglądanie wszystkich rezerwacji w systemie, konfigurowanie reguł anulowania rezerwacji, konfigurowanie wyjątków od reguł konfliktów rezerwacji, generowanie raportów systemowych, zarządzanie rolami i uprawnieniami.
\end{itemize}

% ─── WYMAGANIA FUNKCJONALNE ─────────────────────────────────────────
\section{Wymagania funkcjonalne}

\begin{description}[leftmargin=1.8cm, labelwidth=1.3cm, labelsep=0.5cm, style=nextline]
  \item[FR1] System umożliwia użytkownikom przeglądanie dostępnych terminów wizyt u specjalistów z możliwością filtrowania według specjalizacji, daty i godziny.
  \item[FR2] System umożliwia użytkownikowi dokonanie rezerwacji wybranego terminu u specjalisty, po czym termin ten staje się niedostępny dla innych użytkowników.
  \item[FR3] System zapobiega konfliktom rezerwacji, uniemożliwiając zarezerwowanie tego samego terminu przez więcej niż jednego użytkownika, chyba że administrator skonfigurował wyjątek.
  \item[FR4] Użytkownik może anulować własną rezerwację wyłącznie do 24 godzin przed terminem wizyty. Po upływie tego czasu anulowanie jest zablokowane lub wymaga uzasadnienia.
  \item[FR5] Specjalista może zarządzać swoim grafikiem dostępności, dodając, modyfikując lub usuwając dostępne terminy.
  \item[FR6] System przechowuje historię rezerwacji dla każdego użytkownika i specjalisty, umożliwiając przeglądanie przeszłych wizyt oraz ich statusów (AVAILABLE, BOOKED, CANCELLED, COMPLETED, BLOCKED).
  \item[FR7] Jeden użytkownik może mieć maksymalnie 3 aktywne rezerwacje. Administrator może w uzasadnionych przypadkach zezwolić na przekroczenie tego limitu.
  \item[FR8] Administrator może konfigurować reguły anulowania rezerwacji, w tym minimalny czas przed wizytą wymagany do anulowania oraz warunki dopuszczalnych wyjątków.
  \item[FR9] System powinien obsługiwać równoczesne próby rezerwacji i anulowania, zapewniając możliwie spójny stan danych, jednak szczegóły mechanizmu nie są jednoznacznie określone.
  \item[FR10] System jest zaprojektowany w sposób skalowalny, umożliwiający obsługę rosnącej liczby użytkowników, specjalistów i rezerwacji bez degradacji wydajności.
\end{description}

% ─── WYMAGANIA NIEFUNKCJONALNE ──────────────────────────────────────
\section{Wymagania niefunkcjonalne}

\begin{description}[leftmargin=1.8cm, labelwidth=1.3cm, labelsep=0.5cm, style=nextline]
  \item[NFR1] System powinien obsługiwać co najmniej 10\,000 jednoczesnych użytkowników bez spadku czasu odpowiedzi poniżej 2 sekund dla operacji przeglądania i rezerwacji terminów.
  \item[NFR2] Dostępność systemu powinna wynosić co najmniej 99,9\% w skali miesiąca, z wyłączeniem zaplanowanych okien serwisowych.
  \item[NFR3] Wszystkie dane osobowe muszą być przechowywane zgodnie z RODO, z szyfrowaniem AES-256 w spoczynku i TLS 1.2+ w tranzycie.
  \item[NFR4] System powinien zapewniać spójność stanu końcowego w przypadku równoczesnych operacji rezerwacji tego samego terminu.
  \item[NFR5] Czas odpowiedzi na operację rezerwacji lub anulowania wizyty nie powinien przekraczać 3 sekund w 95. percentylu przy normalnym obciążeniu systemu.
  \item[NFR6] System powinien być zaprojektowany w architekturze umożliwiającej horyzontalne skalowanie komponentów backendowych.
  \item[NFR7] System powinien rejestrować wszystkie operacje krytyczne (rezerwacje, anulowania, modyfikacje grafiku) w niemodyfikowalnym logu audytowym z dokładnością do sekundy.
\end{description}

% ─── USER STORIES ───────────────────────────────────────────────────
\section{User stories}

\begin{description}[leftmargin=1.8cm, labelwidth=1.3cm, labelsep=0.5cm, style=nextline]
  \item[US1] Jako użytkownik chcę przeglądać dostępne terminy wizyt z możliwością filtrowania, aby szybko znaleźć odpowiadający mi termin.
  \item[US2] Jako użytkownik chcę zarezerwować wybrany termin wizyty u specjalisty, aby mieć pewność, że nikt inny nie zajmie tego terminu.
  \item[US3] Jako użytkownik chcę anulować własną rezerwację w dozwolonym oknie czasowym, aby zrezygnować z wizyty bez ponoszenia konsekwencji.
  \item[US4] Jako użytkownik chcę przeglądać historię moich rezerwacji wraz z ich statusami.
  \item[US5] Jako specjalista chcę dodawać, modyfikować i usuwać dostępne terminy w moim grafiku, aby elastycznie zarządzać swoją dostępnością.
  \item[US6] Jako administrator chcę konfigurować wyjątki od reguł konfliktów rezerwacji, aby umożliwić obsługę szczególnych przypadków.
  \item[US7] Jako specjalista chcę przeglądać listę wszystkich moich nadchodzących rezerwacji.
  \item[US8] Jako specjalista chcę blokować wybrane terminy w moim grafiku, aby uniemożliwić rezerwację w czasie, gdy jestem niedostępny.
  \item[US9] Jako administrator chcę konfigurować reguły anulowania rezerwacji, w tym minimalny czas przed wizytą.
  \item[US10] Jako administrator chcę konfigurować wyjątki od reguł konfliktów rezerwacji dla wybranych specjalistów lub terminów.
\end{description}

% ─── ACCEPTANCE CRITERIA ────────────────────────────────────────────
\section{Acceptance criteria}

\textbf{AC1}
\begin{itemize}[itemsep=2pt]
  \item \textbf{Given} użytkownik jest zalogowany i widzi dostępny termin wizyty
  \item \textbf{When} użytkownik klika przycisk rezerwacji wybranego terminu i potwierdza
  \item \textbf{Then} system tworzy rezerwację, zmienia status terminu na BOOKED i wyświetla potwierdzenie
\end{itemize}

\vspace{0.4em}
\textbf{AC2}
\begin{itemize}[itemsep=2pt]
  \item \textbf{Given} dwóch użytkowników jednocześnie próbuje zarezerwować ten sam termin
  \item \textbf{When} obaj użytkownicy klikają przycisk rezerwacji w tym samym czasie
  \item \textbf{Then} system przyznaje rezerwację tylko jednemu użytkownikowi, zapewniając spójność danych
\end{itemize}

\vspace{0.4em}
\textbf{AC3}
\begin{itemize}[itemsep=2pt]
  \item \textbf{Given} użytkownik posiada aktywną rezerwację, a do wizyty pozostało więcej niż 24 godziny
  \item \textbf{When} użytkownik wybiera opcję anulowania rezerwacji
  \item \textbf{Then} system anuluje rezerwację, zwalnia termin (AVAILABLE) i wyświetla potwierdzenie anulowania
\end{itemize}

\vspace{0.4em}
\textbf{AC4}
\begin{itemize}[itemsep=2pt]
  \item \textbf{Given} użytkownik posiada aktywną rezerwację, a do wizyty pozostało mniej niż 24 godziny
  \item \textbf{When} użytkownik próbuje anulować rezerwację
  \item \textbf{Then} system blokuje anulowanie i wyświetla komunikat informujący o przekroczeniu dozwolonego okna czasowego
\end{itemize}

\vspace{0.4em}
\textbf{AC5}
\begin{itemize}[itemsep=2pt]
  \item \textbf{Given} użytkownik posiada już 3 aktywne rezerwacje
  \item \textbf{When} próbuje zarezerwować kolejny termin
  \item \textbf{Then} operacja zostaje odrzucona przez system ze względu na przekroczenie limitu
\end{itemize}

\vspace{0.4em}
\textbf{AC6}
\begin{itemize}[itemsep=2pt]
  \item \textbf{Given} specjalista jest zalogowany i zarządza swoim grafikiem
  \item \textbf{When} specjalista wybiera termin i ustawia blokadę
  \item \textbf{Then} system oznacza wybrany termin jako BLOCKED i uniemożliwia jego rezerwację przez użytkowników
\end{itemize}

\vspace{0.4em}
\textbf{AC7}
\begin{itemize}[itemsep=2pt]
  \item \textbf{Given} specjalista jest zalogowany i znajduje się w panelu zarządzania grafikiem
  \item \textbf{When} specjalista dodaje nowy termin podając datę, godzinę i czas trwania
  \item \textbf{Then} system zapisuje nowy termin jako dostępny i wyświetla go użytkownikom w wynikach wyszukiwania
\end{itemize}

\vspace{0.4em}
\textbf{AC8}
\begin{itemize}[itemsep=2pt]
  \item \textbf{Given} administrator jest zalogowany i znajduje się w panelu konfiguracji
  \item \textbf{When} administrator ustawia minimalny czas przed wizytą wymagany do anulowania
  \item \textbf{Then} system zapisuje nowe reguły i stosuje je do przyszłych rezerwacji
\end{itemize}

\vspace{0.4em}
\textbf{AC9}
\begin{itemize}[itemsep=2pt]
  \item \textbf{Given} administrator jest zalogowany i konfiguruje reguły anulowania
  \item \textbf{When} administrator definiuje wyjątek dla konkretnego specjalisty lub terminu
  \item \textbf{Then} system stosuje wyjątkowe reguły anulowania wyłącznie dla wskazanego przypadku
\end{itemize}

\newpage

% ─── METRYKI VS GOLD ────────────────────────────────────────────────
\section{Metryki porównawcze względem gold artifact (Wpływ trybu Hybrid)}

\begin{mdframed}[backgroundcolor=metabg, linewidth=0pt, innerleftmargin=10pt, innerrightmargin=10pt, innertopmargin=8pt, innerbottommargin=8pt]
\small
\textbf{Wniosek z eksperymentu:} Interwencja ludzka (pruning) drastycznie zwiększyła \textbf{Precision} eliminując halucynacje LLM (tzw. over-engineering), a chirurgiczne wstrzyknięcia kluczowych reguł biznesowych z Gold Standardu podniosły \textbf{Recall}.
\end{mdframed}

\vspace{0.8em}

\subsection{Wymagania funkcjonalne (FR)}

\begin{tabular}{p{2.5cm} p{7cm} l}
\toprule
\textbf{ID Hybrid} & \textbf{Mapowanie na gold} & \textbf{Status} \\
\midrule
FR1 & Gold FR1 -- przeglądanie terminów & \tp{TP} \\
FR2 & Gold FR2 -- rezerwacja dostępnego terminu & \tp{TP} \\
FR3 & Gold FR4 -- brak podwójnej rezerwacji + wyjątki & \tp{TP} \\
FR4 & Gold FR5 -- anulowanie (poprawione na 24h) & \tp{TP} \\
FR5 & Gold FR7 -- zarządzanie grafikiem specjalisty & \tp{TP} \\
FR6 & Gold FR8, FR12 -- statusy i historia (dodano statusy) & \tp{TP} \\
FR7 & Gold FR3 -- limit 3 aktywnych rezerwacji (dodane) & \tp{TP} \\
FR8 & brak w gold -- konfiguracja reguł anulowania przez Admina & \fp{FP} \\
FR9 & Gold FR13 -- równoczesne operacje (dodane) & \tp{TP} \\
FR10 & brak w gold (NFR w gold, nie FR) -- skalowalność & \fp{FP} \\
\midrule
\multicolumn{3}{l}{\textit{Brakujące z gold:}} \\
-- & Gold FR6 -- zwolnienie terminu po anulowaniu & \fn{FN} \\
-- & Gold FR9 -- dostęp do danych według ról & \fn{FN} \\
-- & Gold FR10 -- zakaz nakładających się rezerwacji & \fn{FN} \\
-- & Gold FR11 -- blokady slotów (BLOCKED) explicite & \fn{FN} \\
\bottomrule
\end{tabular}

\vspace{0.5em}
\begin{tabular}{lrrrr}
\toprule
& \textbf{TP} & \textbf{FP} & \textbf{FN} & \\
\midrule
FR & 8 & 2 & 4 & \\
\midrule
\textbf{Precision} & \multicolumn{4}{l}{$8 / (8+2) = \mathbf{0{,}80}$ (wzrost z 0.62)} \\
\textbf{Recall}    & \multicolumn{4}{l}{$8 / (8+4) = \mathbf{0{,}67}$ (wzrost z 0.62)} \\
\textbf{F1}        & \multicolumn{4}{l}{$2 \times 0{,}80 \times 0{,}67 / (0{,}80+0{,}67) = \mathbf{0{,}73}$} \\
\bottomrule
\end{tabular}

\vspace{1em}
\subsection{Wymagania niefunkcjonalne (NFR)}

\begin{tabular}{lrrrr}
\toprule
& \textbf{TP} & \textbf{FP} & \textbf{FN} & \\
\midrule
NFR & 6 & 1 & 0 & \\
\midrule
\textbf{Precision} & \multicolumn{4}{l}{$6 / (6+1) = \mathbf{0{,}86}$ (wzrost z 0.50 po wycięciu szumu)} \\
\textbf{Recall}    & \multicolumn{4}{l}{$6 / (6+0) = \mathbf{1{,}00}$ (wzrost po dodaniu obsługi współbieżności)} \\
\textbf{F1}        & \multicolumn{4}{l}{$2 \times 0{,}86 \times 1{,}00 / (0{,}86+1{,}00) = \mathbf{0{,}92}$} \\
\bottomrule
\end{tabular}

\vspace{1em}
\subsection{Acceptance criteria (AC)}

\begin{tabular}{p{2.5cm} p{7cm} l}
\toprule
\textbf{ID Hybrid} & \textbf{Mapowanie na gold} & \textbf{Status} \\
\midrule
AC1 & Gold AC1 -- rezerwacja $\rightarrow$ BOOKED & \tp{TP} \\
AC2 & Gold AC3 -- double booking spójność & \tp{TP} \\
AC3--AC4 & Gold AC4 -- anulowanie $>$24h oraz $<$24h & \tp{TP} \\
AC5 & Gold AC2 -- limit 3 rezerwacji (dodane) & \tp{TP} \\
AC6--AC9 & brak w gold -- dodatkowe przypadki dla Specjalisty/Admina & \fp{FP} \\
\midrule
-- & Gold AC5 -- anulowanie graniczne (dokładnie 24h) & \fn{FN} \\
-- & Gold AC6 -- konflikt czasowy nakładających się rezerwacji & \fn{FN} \\
\bottomrule
\end{tabular}

\vspace{0.5em}
\begin{tabular}{lrrrr}
\toprule
& \textbf{TP} & \textbf{FP} & \textbf{FN} & \\
\midrule
AC & 4 & 5 & 2 & \\
\midrule
\textbf{Precision} & \multicolumn{4}{l}{$4 / (4+5) = \mathbf{0{,}44}$ (znaczny wzrost z 0.13 po redukcji z 23 do 9 AC)} \\
\textbf{Recall}    & \multicolumn{4}{l}{$4 / (4+2) = \mathbf{0{,}67}$ (wzrost po uzupełnieniu logiki limitów i 24h)} \\
\textbf{F1}        & \multicolumn{4}{l}{$2 \times 0{,}44 \times 0{,}67 / (0{,}44+0{,}67) = \mathbf{0{,}53}$} \\
\bottomrule
\end{tabular}

% ─── METADANE I OCENA ───────────────────────────────────────────────
\section{Metadane i ocena}
\begin{table}[h]
\begin{tabularx}{\textwidth}{l X}
\toprule
\textbf{Pole} & \textbf{Wartość} \\
\midrule
Model & \texttt{claude-sonnet-4-6} \\
Tryb & \textbf{Hybrid} (LLM generacja + Human-in-the-loop pruning) \\
Case study & appointment-booking (ver2 challenging) \\
\midrule
\multicolumn{2}{l}{\textit{Ocena (0--3)}} \\
\midrule
Correctness (M1)     & \scorebox{3} \quad błędy biznesowe (24h, limit 3) zostały załatane przez człowieka \\
Completeness (M2)    & \scorebox{3} \quad wymogi architektoniczne pokryte znacznie lepiej (współbieżność, statusy) \\
Consistency (M3)     & \scorebox{3} \quad w pełni spójny \\
Clarity (M4)         & \scorebox{3} \quad zachowany balans - brak "ściany tekstu" po redukcji halucynacji LLM \\
Maintainability (M5) & \scorebox{3} \quad artefakt gotowy do przekazania do Architecture Team \\
\bottomrule
\end{tabularx}
\end{table}

\end{document}